%% ma: L12-B08-0109
%% lop: 12
%% chuong: 2
%% bai: 8
%% dang: 12.8.27
%% loai: TN4
%% muc: TH
%% dap_an: "A"
%% nguon: "CĐ2.3.pdf giáo viên gửi 27/10/2026 (M07, Lê Minh Khôi – Nguyễn Anh Khôi, Chương II Chủ đề 3 Biểu thức tọa độ của các phép toán vectơ), Câu 20"
%% kien_thuc: "Bài 8 (tọa độ vectơ), Bài 6 (tính chất trọng tâm tam giác)"
%% trang_thai: chinh-thuc
\begin{ex}
Trong không gian $Oxyz$, cho tam giác $ABC$ có đỉnh $A(2;4;-3)$ và $G(2;1;0)$ là trọng tâm của tam giác. Tọa độ của vectơ $\vec u=\overrightarrow{AB}+\overrightarrow{AC}$ là
\choice{\True $\vec u=(0;-9;9)$}{$\vec u=(0;-4;4)$}{$\vec u=(0;4;-4)$}{$\vec u=(0;9;-9)$}
\loigiai{$\overrightarrow{AB}+\overrightarrow{AC}=3\overrightarrow{AG}=3(0;-3;3)=(0;-9;9)$.}
\end{ex}
